% GENERATED FILE. Edit canonical lessons, then run npm run generate:books.
% canonical-source-hash: fnv1a64:1893feb535b5b225
% canonical-lessons: TA-C350-kannottam, TA-C350-nilaippatu, TA-C350-manappanmai, TA-C350-konam, TA-C350-velippatuttu

\chapter{Points of View}
\label{ch:ta-points-of-view}
\hlchaptermodality{350}

\begin{chapteropening}
\textbf{By the end of this chapter:} \emph{I can talk about a point of view, a stand and an attitude.}

Complete the last lesson of chapter 350: I can talk about a point of view, a stand and an attitude.
\end{chapteropening}

\section[\texorpdfstring{\ta{கண்ணோட்டம்} (kaṇṇōṭṭam)}{கண்ணோட்டம் (kaṇṇōṭṭam)}]{\ta{கண்ணோட்டம்} (kaṇṇōṭṭam) — a point of view; a perspective}

\label{lesson:TA-C350-kannottam}



\begin{quote}
Before the new one: say the Tamil for a hypothesis, then the Tamil for many people.
\end{quote}

\subsection*{You'll want to know: \ta{கண்ணோட்டம்}}
\textbf{\ta{கண்ணோட்டம்}} — \emph{kaṇṇōṭṭam} — \textquotedblleft{}a point of view; a perspective\textquotedblright{}.

\begin{quote}
\ta{அவருடைய} \ta{கண்ணோட்டம்} \ta{எனக்குப்} \ta{புரிகிறது}.
\end{quote}

\begin{quote}
\emph{avaruṭaiya kaṇṇōṭṭam eṉakkup purikiṟatu.}
\end{quote}

\begin{quote}
I understand his point of view.
\end{quote}

\ta{வேறு} \ta{கண்ணோட்டத்தில்}, from another point of view.

\begin{cousinweb}
\textbf{\ta{கண்ணோட்டம்}} (\emph{kaṇṇōṭṭam}) is a native compound of \emph{kaṇ}, \textquotedblleft{}eye\textquotedblright{} (DEDR 1159), and \emph{ōṭṭam}, \textquotedblleft{}running, course\textquotedblright{} (DEDR 1041): where the eye runs.
\end{cousinweb}

\begin{grammarlens}[title={What you've built}]
One more word for giving reasons and opinions.
\end{grammarlens}

\subsection*{Your turn}
\begin{itemize}
\raggedright
  \item \emph{Say it:} \emph{kaṇṇōṭṭam}
  \item \emph{Say it:} \emph{kaṇṇōṭṭam}, once more
  \item \emph{Say it:} \emph{avaruṭaiya kaṇṇōṭṭam eṉakkup purikiṟatu.} — I understand his point of view.
  \item \emph{Recall:} say the Tamil for a hypothesis, then the Tamil for many people, then say \emph{kaṇṇōṭṭam} again
\end{itemize}

\begin{tcolorbox}[breakable,colback=teal!4,colframe=teal!35!black,title={Before you move on}]
What does \ta{கண்ணோட்டம்} mean? (\textquotedblleft{}A point of view; a perspective\textquotedblright{}.) Say it once more.
\end{tcolorbox}
\section[\texorpdfstring{\ta{நிலைப்பாடு} (nilaippāṭu)}{நிலைப்பாடு (nilaippāṭu)}]{\ta{நிலைப்பாடு} (nilaippāṭu) — a stand; a position}

\label{lesson:TA-C350-nilaippatu}



\begin{quote}
Before the new one: say the Tamil for many people, then the Tamil for a point of view.
\end{quote}

\subsection*{You'll want to know: \ta{நிலைப்பாடு}}
\textbf{\ta{நிலைப்பாடு}} — \emph{nilaippāṭu} — \textquotedblleft{}a stand; a position\textquotedblright{}.

\begin{quote}
\ta{இந்த} \ta{விஷயத்தில்} \ta{உங்கள்} \ta{நிலைப்பாடு} \ta{என்ன}?
\end{quote}

\begin{quote}
\emph{inta viṣayattil uṅkaḷ nilaippāṭu eṉṉa?}
\end{quote}

\begin{quote}
What is your stand on this matter?
\end{quote}

\ta{நிலைப்பாட்டை} \ta{மாற்று}, to change one's stand: a -\ta{டு} noun doubles before a case ending.

\begin{cousinweb}
\textbf{\ta{நிலைப்பாடு}} (\emph{nilaippāṭu}) is native: \emph{nilai}, \textquotedblleft{}a standing, a state\textquotedblright{}, belongs to \emph{nil}, \textquotedblleft{}to stand\textquotedblright{} (DEDR 3675), and \emph{pāṭu} is the noun of \emph{paṭu}: \textquotedblleft{}the place where one stands\textquotedblright{}.
\end{cousinweb}

\begin{grammarlens}[title={What you've built}]
One more word for giving reasons and opinions.
\end{grammarlens}

\subsection*{Your turn}
\begin{itemize}
\raggedright
  \item \emph{Say it:} \emph{nilaippāṭu}
  \item \emph{Say it:} \emph{nilaippāṭu}, once more
  \item \emph{Say it:} \emph{inta viṣayattil uṅkaḷ nilaippāṭu eṉṉa?} — What is your stand on this matter?
  \item \emph{Recall:} say the Tamil for many people, then the Tamil for a point of view, then say \emph{nilaippāṭu} again
\end{itemize}

\begin{tcolorbox}[breakable,colback=teal!4,colframe=teal!35!black,title={Before you move on}]
What does \ta{நிலைப்பாடு} mean? (\textquotedblleft{}A stand; a position\textquotedblright{}.) Say it once more.
\end{tcolorbox}
\section[\texorpdfstring{\ta{மனப்பான்மை} (maṉappāṉmai)}{மனப்பான்மை (maṉappāṉmai)}]{\ta{மனப்பான்மை} (maṉappāṉmai) — an attitude}

\label{lesson:TA-C350-manappanmai}



\begin{quote}
Before the new one: say the Tamil for a point of view, then the Tamil for a stand.
\end{quote}

\subsection*{You'll want to know: \ta{மனப்பான்மை}}
\textbf{\ta{மனப்பான்மை}} — \emph{maṉappāṉmai} — \textquotedblleft{}an attitude\textquotedblright{}.

\begin{quote}
\ta{அவனுடைய} \ta{மனப்பான்மை} \ta{சரியில்லை}.
\end{quote}

\begin{quote}
\emph{avaṉuṭaiya maṉappāṉmai cariyillai.}
\end{quote}

\begin{quote}
His attitude is not right.
\end{quote}

\ta{நல்ல} \ta{மனப்பான்மை}, a good attitude.

\begin{cousinweb}
\textbf{\ta{மனப்பான்மை}} (\emph{maṉappāṉmai}) is a compound: \emph{maṉam}, \textquotedblleft{}mind\textquotedblright{}, is borrowed from Sanskrit \emph{manas} (CDIAL 9822), and native \emph{pāṉmai}, \textquotedblleft{}portion; nature\textquotedblright{} (DEDR 4097), gives \textquotedblleft{}the nature of the mind\textquotedblright{}.
\end{cousinweb}

\begin{grammarlens}[title={What you've built}]
One more word for giving reasons and opinions.
\end{grammarlens}

\subsection*{Your turn}
\begin{itemize}
\raggedright
  \item \emph{Say it:} \emph{maṉappāṉmai}
  \item \emph{Say it:} \emph{maṉappāṉmai}, once more
  \item \emph{Say it:} \emph{avaṉuṭaiya maṉappāṉmai cariyillai.} — His attitude is not right.
  \item \emph{Recall:} say the Tamil for a point of view, then the Tamil for a stand, then say \emph{maṉappāṉmai} again
\end{itemize}

\begin{tcolorbox}[breakable,colback=teal!4,colframe=teal!35!black,title={Before you move on}]
What does \ta{மனப்பான்மை} mean? (\textquotedblleft{}An attitude\textquotedblright{}.) Say it once more.
\end{tcolorbox}
\section[\texorpdfstring{\ta{கோணம்} (kōṇam)}{கோணம் (kōṇam)}]{\ta{கோணம்} (kōṇam) — an angle}

\label{lesson:TA-C350-konam}



\begin{quote}
Before the new one: say the Tamil for a stand, then the Tamil for an attitude.
\end{quote}

\subsection*{You'll want to know: \ta{கோணம்}}
\textbf{\ta{கோணம்}} — \emph{kōṇam} — \textquotedblleft{}an angle\textquotedblright{}.

\begin{quote}
\ta{அவர்} \ta{இந்தப்} \ta{பிரச்சினையை} \ta{வேறு} \ta{கோணத்தில்} \ta{பார்த்தார்}.
\end{quote}

\begin{quote}
\emph{avar intap piracciṉaiyai vēṟu kōṇattil pārttār.}
\end{quote}

\begin{quote}
He saw this problem from a different angle.
\end{quote}

The same word is an angle in geometry; \ta{வேறு} \ta{கோணத்தில்}, from another angle.

\begin{cousinweb}
\textbf{\ta{கோணம்}} (\emph{kōṇam}) is borrowed from Sanskrit \emph{koṇa}, \textquotedblleft{}corner, angle\textquotedblright{} (CDIAL 3504). CDIAL itself notes that the Sanskrit word may have come from Dravidian, so the word may have gone north and come back.
\end{cousinweb}

\begin{grammarlens}[title={What you've built}]
One more word for giving reasons and opinions.
\end{grammarlens}

\subsection*{Your turn}
\begin{itemize}
\raggedright
  \item \emph{Say it:} \emph{kōṇam}
  \item \emph{Say it:} \emph{kōṇam}, once more
  \item \emph{Say it:} \emph{avar intap piracciṉaiyai vēṟu kōṇattil pārttār.} — He saw this problem from a different angle.
  \item \emph{Recall:} say the Tamil for a stand, then the Tamil for an attitude, then say \emph{kōṇam} again
\end{itemize}

\begin{tcolorbox}[breakable,colback=teal!4,colframe=teal!35!black,title={Before you move on}]
What does \ta{கோணம்} mean? (\textquotedblleft{}An angle\textquotedblright{}.) Say it once more.
\end{tcolorbox}
\section[\texorpdfstring{\ta{வெளிப்படுத்து} (veḷippaṭuttu)}{வெளிப்படுத்து (veḷippaṭuttu)}]{\ta{வெளிப்படுத்து} (veḷippaṭuttu) — to express; to reveal}

\label{lesson:TA-C350-velippatuttu}



\begin{quote}
Before the new one: say the Tamil for an attitude, then the Tamil for an angle.
\end{quote}

\subsection*{You'll want to know: \ta{வெளிப்படுத்து}}
\textbf{\ta{வெளிப்படுத்து}} — \emph{veḷippaṭuttu} — \textquotedblleft{}to express; to reveal\textquotedblright{}.

\begin{quote}
\ta{அவள்} \ta{தன்} \ta{கருத்தைத்} \ta{தைரியமாக} \ta{வெளிப்படுத்தினாள்}.
\end{quote}

\begin{quote}
\emph{avaḷ taṉ karuttait tairiyamāka veḷippaṭuttiṉāḷ.}
\end{quote}

\begin{quote}
She expressed her opinion boldly.
\end{quote}

Past (I …): \ta{வெளிப்படுத்தினேன்}, \emph{veḷippaṭuttiṉēṉ}.

What is expressed takes -\ta{ஐ}: \ta{கருத்தை} \ta{வெளிப்படுத்து}.

\begin{cousinweb}
\textbf{\ta{வெளிப்படுத்து}} (\emph{veḷippaṭuttu}) is native: \emph{veḷi}, \textquotedblleft{}outside, open space\textquotedblright{} (DEDR 5498), with \emph{paṭuttu}, the causative of \emph{paṭu}: \textquotedblleft{}to bring out into the open\textquotedblright{}.
\end{cousinweb}

\begin{grammarlens}[title={What you've built}]
One more word for giving reasons and opinions.
\end{grammarlens}

\subsection*{Your turn}
\begin{itemize}
\raggedright
  \item \emph{Say it:} \emph{veḷippaṭuttu}
  \item \emph{Say it:} \emph{veḷippaṭuttu}, once more
  \item \emph{Say it:} \emph{avaḷ taṉ karuttait tairiyamāka veḷippaṭuttiṉāḷ.} — She expressed her opinion boldly.
  \item \emph{Recall:} say the Tamil for an attitude, then the Tamil for an angle, then say \emph{veḷippaṭuttu} again
\end{itemize}

\begin{tcolorbox}[breakable,colback=teal!4,colframe=teal!35!black,title={Before you move on}]
What does \ta{வெளிப்படுத்து} mean? (\textquotedblleft{}To express; to reveal\textquotedblright{}.) Say it once more.
\end{tcolorbox}
